% Generated by scripts/txt2tex.py from chapters/abstract.txt. Edit the TXT, not this file.
הספר מתחיל משאלה שכל אחד יכול לשאול: מישהו מקליד למערכת שאלה על בצק שלא תופח, ומקבל תשובה בטוחה ומסודרת. מה קרה בין ההקלדה לתשובה? הקוראת עוברת בשבעה פרקים מאחורי המסך, צעד אחר צעד, ובכל פרק נחשף חלק אחר במנגנון: מה זה בכלל מכונה שלומדת, איך נאספים נתונים ומה קורה כשהם מוטים, איך רשת נוירונים מתאימה את עצמה, איך מודל שפה מנבא מילה אחר מילה, ואיך בודקים אם מכונה באמת למדה ולא רק נשמעת כך.

הפרק על הבדיקה כולל ניסוי קטן של הספר עצמו, שבו המספרים מוצגים כפי שיצאו. אחריו באים הגבולות: מה קורה למה שהקוראת הקלידה, מתי מערכת טועה דווקא בטון הבטוח ביותר, ומה כדאי לדעת לפני שסומכים עליה. הפרק האחרון חוזר אל המסך ומציע לא כללים לשינון אלא שיקול דעת: מתי לסמוך, מתי לבדוק, ואיך להשתמש בכלי בלי לוותר על החשיבה של הקוראת עצמה.

נספח נפרד מספר איך נוצר הספר עצמו: בעזרת סוכנים של בינה מלאכותית, שכל אחד מהם עשה תפקיד אחד, תחת פיקוח אנושי ועם תיעוד פתוח לבדיקה. כך הספר מראה בעצמו את מה שהוא מלמד, שאמון נבנה מבדיקה ולא מביטחון.
